\section{Preliminaries}
\label{sec:prelim}

We introduce the notation and building blocks used throughout the paper.  Let
$\mathbb{F}$ be a finite field of prime order $p$, let $\lambda$ be the
security parameter, and write $\textsf{negl}(\lambda)$ and $\ppt$ for a
negligible function and probabilistic polynomial time, respectively.

\subsection{Notation}
\label{sec:notation}

We use bold uppercase letters for matrices and bold lowercase letters for
vectors, which are row vectors unless stated otherwise.  For a matrix
$\mat{M}$, $\mat{M}[i,j]$, $\mat{M}[i,*]$, and $\mat{M}[*,j]$ denote an entry,
a row, and a column, respectively.  We write $[n]=\{0,1,\ldots,n-1\}$;
$I\sample[n]^t$ denotes $t$ independently sampled indices from $[n]$, and
$r\sample\mathbb{F}$ denotes a uniform field element.

\subsection{Linear Error-Correcting Codes}
\label{sec:prelim-codes}

A linear $[n,k]$ code $\mathcal{C}\subseteq\mathbb{F}^n$ is the image of a
linear encoder $\enc:\mathbb{F}^k\to\mathbb{F}^n$.  Its rate is
$\rho=k/n$.  Let $d_H$ denote Hamming distance and
$\Delta(\vect{u},\vect{v})=d_H(\vect{u},\vect{v})/n$ relative Hamming
distance.  If $d_{\min}$ is the minimum Hamming distance of $\mathcal{C}$,
then its relative distance and unique-decoding radius are
$\delta=d_{\min}/n$ and
$\tau_\ud=n^{-1}\lfloor(d_{\min}-1)/2\rfloor$, respectively.  We use a
proximity threshold $\tau\leq\tau_\ud$.

For $\mat{M}\in\mathbb{F}^{k\times k}$, its row-wise encoding
$\widehat{\mat{M}}=\enc(\mat{M})\in\mathbb{F}^{k\times n}$ has rows
$\enc(\mat{M}[i,*])$ for $i\in[k]$.

\begin{definition}[Interleaved code]\label{def:interleaved-code}
For $\ell\geq1$, the $\ell$-fold interleaved code
$\mathcal{C}^{\langle\ell\rangle}$ groups $\ell$ codewords position-wise:
its $j$-th symbol is
$(\vect{c}^{(0)}[j],\ldots,\vect{c}^{(\ell-1)}[j])\in\mathbb{F}^{\ell}$.
Thus $\widehat{\mat{M}}$ is a $k$-fold interleaved codeword whose $j$-th
symbol is $\widehat{\mat{M}}[*,j]$.
\end{definition}

For $\vect{u}\in\mathbb{F}^k$, linearity gives
\[
  \enc(\vect{u}\mat{M})[j]
  =\sum_{i\in[k]}u_i\widehat{\mat{M}}[i,j]
  \eqqcolon\fold(\vect{u},\widehat{\mat{M}}[*,j]).
\]
Our construction applies to linear codes satisfying the required proximity-gap
property.  The implementation uses an $[n,k]$ Reed--Solomon code with relative
distance $\delta_{\mathsf{RS}}=(n-k+1)/n$.

\subsection{Commitment Schemes}
\label{sec:prelim-commitments}

A commitment scheme binds a committer to a message while optionally hiding it.
We use Merkle trees for position binding and Pedersen commitments for
computational binding, perfect hiding, and additive homomorphism.

\paragraph{Merkle trees.}
For a collision-resistant hash function $H$, we write
$\rt\gets\merkleC((v_j)_{j\in[n]};\bot)$ for a Merkle commitment with
deterministic leaves, $\pi_j\gets\merkleO((v_i)_{i\in[n]},j)$ for an
authentication path, and $\merkleV(\rt,j,v_j,\pi_j)$ for verification.  Under
collision resistance, a fixed root is position-binding: it cannot be opened to
two different values at the same index except with negligible probability.

\paragraph{Pedersen commitments.}
Let $\mathbb{G}$ be a group of order $|\mathbb{F}|$ and
$\ck=(G_0,\ldots,G_{\ell-1},H)$ a commitment key.  For
$\vect{v}\in\mathbb{F}^{\ell}$ and $o\sample\mathbb{F}$, define
\[
  \pedC(\ck,\vect{v};o)
  =\sum_{i\in[\ell]}\vect{v}[i]G_i+oH.
\]
Pedersen commitments are perfectly hiding, computationally binding under the
discrete-logarithm assumption, and additively homomorphic.


\subsection{Succinct Non-interactive Arguments of Knowledge}
\label{sec:prelim-snark}

A SNARK (succinct non-interactive argument of knowledge) for an NP relation
$\relation$ is a tuple of algorithms $\bPi=(\setup,\prove,\verify)$.

\begin{itemize}
\item $\para \gets \bPi.\setup(1^\lambda,\relation)$:
generates public parameters $\para$ for the relation $\relation$.

\item $\pi \gets \bPi.\prove(\para,\stmt,\wit)$:
outputs a proof $\pi$ for a statement $\stmt$ and witness $\wit$ satisfying $(\stmt,\wit)\in\relation$.

\item $b \gets \bPi.\verify(\para,\stmt,\pi)$:
outputs $b=1$ if $\pi$ is accepted as a valid proof for $\stmt$, and outputs $b=0$ otherwise.
\end{itemize}

\begin{remark}[Security Properties]
We use the standard SNARK notions of completeness, knowledge soundness, and
zero knowledge.  Appendix~\ref{app:snark-security} recalls these notions
formally, including their negligible error bounds. 
\end{remark}


\subsubsection{Commit-and-Prove SNARKs}
A commit-and-prove SNARK (CP-SNARK)~\cite{gnark, legosnark} is a variant of a SNARK that proves statements about witnesses bound to an external commitment. We call such a commitment a \emph{witness commitment}. Let $\com$ be a commitment scheme with commitment key $\ck$. For an NP relation $\relation$ over tuples $(\stmt,\comwit,\omega)$, a CP-SNARK proves knowledge of a witness
$\wit := (\comwit,\omega)$ such that $(\stmt,\comwit,\omega)\in\relation$,
where $\comwit$ is the committed part of the witness and $\omega$ is the
remaining private witness.  The committed witness is bound separately by a
witness commitment
$\widetilde{\cm} = \com(\ck_S, m_0 || ... || m_{q-1}; \widetilde{o})$, where $\tilde{o}$ is the
opening randomness.  Equivalently, the CP proof is verified with respect to the
extended public statement $\stmt_{\cp}:=(\stmt,\widetilde{\cm})$.  When the
context is clear, we write this extended statement simply as $\stmt$.
A CP-SNARK scheme is a tuple of algorithms
$\bPi_{\cp}=(\setup,\com,\prove,\verify)$.

\begin{itemize}
    \item $\para \gets \bPi_{\cp}.\setup(1^\lambda,\relation)$: generates public parameters $\para=(\pk,\vk,\ck)$ for the relation $\relation$, where $\pk$ is the proving key, $\vk$ is the verifying key, and $\ck$ is the commitment key.

    \item $\widetilde{\cm} \gets \bPi_{\cp}.\com(\para,\comwit;\tilde{o})$:
    outputs a witness commitment to the committed witness part $\comwit$.

    \item $\pi_{\cp} \gets \bPi_{\cp}.\prove(\para,\stmt_{\cp},\wit)$:
    outputs a proof for the extended statement $\stmt_{\cp}$ and witness
    $\wit=(\comwit,\omega)$.

    \item $b \gets \bPi_{\cp}.\verify(\para,\stmt_{\cp},\pi_{\cp})$:
    outputs $b=1$ if $\pi_{\cp}$ is accepted as a valid proof for the extended
    statement $\stmt_{\cp}$, and outputs $b=0$ otherwise. Equivalently, we may
    write verification as
    $\bPi_{\cp}.\verify(\para,\stmt,\widetilde{\cm},\pi_{\cp})$.
\end{itemize}
\subsection{CP-Link}
\label{sec:prelim-cplink}

A CP-Link proof is a proof system for linking commitments.  In the form used by
\(\protocol\), it proves that one SNARK-side commitment to an ordered
concatenation of sampled blocks is consistent with external commitments to the
individual blocks, without revealing the blocks or the opening randomness.

Let $\com$ be a commitment scheme.  For a SNARK-side key \(\ck_S\), an external
key \(\ck_E\), a SNARK-side commitment \(\widetilde{\cm}\), and external
commitments \(\cm_0,\ldots,\cm_{q-1}\), the blockwise linking relation is
defined over the statement and witness
\[
\begin{aligned}
\stmt_{\mathsf{link}}
&=
((\ck_S,\widetilde{\cm}),
(\ck_E,\cm_0,\ldots,\cm_{q-1})),\\
\wit_{\mathsf{link}}
&=
((\vect{m}_i)_{i=0}^{q-1},\widetilde{o},(o_i)_{i=0}^{q-1}).
\end{aligned}
\]
The pair \((\stmt_{\mathsf{link}},\wit_{\mathsf{link}})\) is in
\(\relation_{\mathsf{link}}\) if and only if
\[
\begin{aligned}
\widetilde{\cm}
&=
\com(\ck_S,\vect{m}_0\|\cdots\|\vect{m}_{q-1};\widetilde{o}),\\
\cm_i
&=
\com(\ck_E,\vect{m}_i;o_i)
\quad \text{for all } i=0,\ldots,q-1.
\end{aligned}
\]
A proof for $\relation_{\mathsf{link}}$ convinces the verifier that the ordered
blocks committed inside the CP-SNARK are exactly the blocks committed by the
corresponding external commitments.

We use CP-Link through the standard statement/witness interface for a linking
relation.

A CP-Link proof system is a tuple of algorithms
$\bPi_{\mathsf{link}}=(\setup,\prove,\verify)$.

\begin{itemize}
\item $\para_{\mathsf{link}} \gets
\bPi_{\mathsf{link}}.\setup(1^\lambda,\relation_{\mathsf{link}})$:
generates public parameters for the linking relation.

\item $\pi_{\mathsf{link}} \gets
\bPi_{\mathsf{link}}.\prove(
\para_{\mathsf{link}},
\stmt_{\mathsf{link}},
\wit_{\mathsf{link}})$:
outputs a proof that the statement-witness pair satisfies the linking relation.

\item $b \gets
\bPi_{\mathsf{link}}.\verify(
\para_{\mathsf{link}},
\stmt_{\mathsf{link}},
\pi_{\mathsf{link}})$:
outputs $b=1$ if the proof is accepted, and outputs $b=0$ otherwise.
\end{itemize}

\begin{remark}[Security Properties]
For the soundness of \(\protocol\), we require CP-Link soundness: an accepting
linking proof should imply membership in \(\relation_{\mathsf{link}}\), except
with error \(\epsilon_{\mathsf{link}}(\lambda)\).  If the linking backend is used
to protect the linked blocks and openings, one may additionally require witness
zero knowledge.  Appendix~\ref{app:cplink-security} states these properties.
\end{remark}

\subsection{Freivalds' Algorithm}
\label{sec:prelim-freivalds}

Freivalds' algorithm~\cite{freivalds1977probabilistic} is a classical randomized algorithm for verifying matrix multiplication.
Given $\mat{A}, \mat{B}, \mat{C} \in \mathbb{F}^{k \times k}$, it samples $\vect{r} \sample \mathbb{F}^k$ and checks whether $\vect{r}(\mat{A}\mat{B}) = \vect{r}\mat{C}$.
The correctness relies on the following:

\begin{lemma}[Schwartz--Zippel~\cite{schwartz1980fast,zippel1979probabilistic}]\label{lem:sz}
Let $P(X_1, \ldots, X_m)$ be a non-zero polynomial of total degree $d$ over a finite field $\mathbb{F}$.
Then for $r_1, \ldots, r_m \sample \mathbb{F}$ chosen uniformly and independently,
\[
\Pr[P(r_1, \ldots, r_m) = 0] \le \frac{d}{|\mathbb{F}|}.
\]
\end{lemma}

In our protocol, we use a structured variant: instead of sampling $\vect{r} \sample \mathbb{F}^k$ uniformly, we sample one scalar $r \sample \mathbb{F}$ and set $\vect{r} = (1, r, r^2, \ldots, r^{k-1})$.
This keeps the Freivalds challenge compact in the transcript and in the SNARK
statement: the verifier records one field element rather than $k$ independent
ones. If $\mat{D} = \mat{A}\mat{B} - \mat{C} \neq \mat{0}$, then for any non-zero
column $\vect{d}_j$ of $\mat{D}$, the value $\langle \vect{r}, \vect{d}_j \rangle$ is a non-zero univariate polynomial in
$r$ of degree at most $k-1$.
By Lemma~\ref{lem:sz}, it vanishes with probability at most $(k-1)/|\mathbb{F}|$.
